\documentclass[12pt]{article}

\usepackage[margin = .8in]{geometry}
\usepackage{amsmath}
\usepackage{graphicx}
\usepackage{multicol, enumerate, tabularx}

\usepackage{adjustbox, soul, setspace}

\usepackage{fancyhdr}
\pagestyle{fancy}

\renewcommand{\emph}[1]{\textsf{\textbf{#1}}} %%% bold emph


\lhead{Math F113X: Numbers and Society}
\rhead{Date: \hspace{1in}}

\usepackage{tikz}
\usetikzlibrary{calc,trees,positioning,arrows,fit,shapes,through, backgrounds}
\usetikzlibrary{patterns}

\usetikzlibrary{decorations.markings}
\usetikzlibrary{arrows}

\usepackage{pgfplots}

\usepackage{longtable}
\usepackage{tabularx}

\newcommand{\ds}{\displaystyle}
\newcommand{\ans}[1][1in]{\rule{#1}{.5pt}}

\newcommand{\points}[1]{(#1 points.)}		% Trying to be lazy.

\usepackage{array}
\newcolumntype{L}[1]{>{\raggedright\let\newline\\\arraybackslash\hspace{0pt}}m{#1}}
\newcolumntype{C}[1]{>{\centering\let\newline\\\arraybackslash\hspace{0pt}}m{#1}}
\newcolumntype{R}[1]{>{\raggedleft\let\newline\\\arraybackslash\hspace{0pt}}m{#1}}
\newcommand{\red}[1]{\textcolor{red}{#1}}

\newcommand{\be}{\begin{enumerate}}
\newcommand{\ee}{\end{enumerate}}

\begin{document}
\begin{center}
{\Large  Worksheet Finance 4: Summary of Topics}
\end{center}

\begin{enumerate}
\item Below are the formulas you will be given on Exam 3. For each scenario below,  \emph{plug in numbers into the appropriate formula} needed to calculate the value. Do \emph{not} actually calculate the the value.

\begin{center}
\textbf{Formulas}
$$ A=P+I \quad \hspace{1cm} \quad A=P(1+rt) \quad \hspace{1cm} \quad A=P\left(1+\frac{r}{n}\right)^{(nt)} \quad \hspace{1cm} \quad P=\frac{A}{\left( 1+\frac{r}{n}\right)^{(nt)}} $$

\vspace*{-.5in}
$$ P=\frac{d\left(1-\left(1+\frac{r}{n}\right)^{(-nt)}\right)}{\left(\frac{r}{n} \right)} \hspace{1cm} \quad d= \frac{P\left(\frac{r}{n} \right)}{\left( 1- \left(1+\frac{r}{n}\right)^{(-nt)}\right)}$$
\end{center}
	\begin{enumerate}
	\item You want to take out a loan to buy a \$150,000 home. The bank offers a 30-year mortgage with an interest rate of 5.2\% compounded monthly. What will the monthly payments be?
	\vfill
	\item You loan your friend \$900. They agree to pay an annual interest rate of 4\% \emph{simple interest}. Eighteen months later, they repay the loan. How much did they pay you in total?
	\vfill
	\item You deposit \$2000 in a savings account with an APR of 3.2\%  compounded quarterly. How much will the account be worth in 10 years?
	\vfill
	\item Suppose an investment typically earns 10\% interest compounded yearly. How much would your initial investment ( or \emph{principal}) need to be in order for the future value of the account to be worth \$100,000 at the end of 40 years?
	\vfill	 
	\end{enumerate}
\newpage
\item For the questions below, \emph{write out the appropriate calculation} and use a calculator to find the numerical answer.
	\begin{enumerate}
	\item Someone tells you that 2 years after depositing \$1200 in an account, the value of the account increased by 28\%. How much money was added to the account over this 2-year period?
	\vfill
	\item A savings account earns 3\% interest each year. If someone initially invests \$1400, what is the value of the account in 30 years?
	\vfill
	\item If you have a mortgage and you pay exactly \$1200 a month for 30 years, how much did you pay over the life of the mortgage? 
	 \vfill
	\end{enumerate}
\item These problems should be completed using a spreadsheet.
	\begin{enumerate}
	\item Suppose someone has a balance of \$2000 on their credit card which has an APR of 27\% compounded monthly. Suppose the person pays the \emph{minimum payment} which is typically interest plus 3\% of the remaining balance. (In fact, the minimum payment usually has a lower bound of \$25 or similar.}			\begin{enumerate}
		\item How long will it take till the balance on the credit card drops below \$1000 purchase assuming no additional purchases are added to the credit card? 
		\vfill
		\item How much in total will have been paid to the credit card company at this point?
		\vfill
		\end{enumerate}
	\item (Bonus) Suppose you open a retirement account with an initial deposit of \$1000 and each month you add \$100. If this account earns 4\% interest compounded monthly, how much money will be in the account in 30 years assuming there are no withdrawals?
	\vfill
	\end{enumerate}

\end{enumerate}
\end{document}

